\scanpage{24}
\mspage{23}

\uline{Proposition} Soient $F,G\in\unsure{\mathcal F|\mathfrak F|\mathscr F}$. Conditions équivalentes
\[
\begin{aligned}
&\text{(i)}\quad F\leq G\\
&\text{(i bis)}\quad \unsure{\widetilde F|\widetilde{\widetilde F}}\subset\widetilde G\\
&\text{(ii)}\quad \mathrm{Mul\unsure{t|\deleted{t}}omb}(F)\leq\mathrm{Mul\unsure{t|\deleted{t}}omb}(G)
 \qquad\begin{gathered}\text{i.e. }\mathrm{Multomb}(F)\\\text{est une sous-figure}\\\text{de }\mathrm{Multomb}(G)\end{gathered}\\
&\text{(ii bis)}\quad \mathrm{Mul\unsure{t|\deleted{t}}omb}(F)\subset\mathrm{Mul\unsure{t|\deleted{t}}omb}\unsure{G)|(G)}
\end{aligned}
\]

Dém. $F\leq G$ \deleted{trivial, et} $\unsure{\Longleftrightarrow|\Longleftarrow}\widetilde F\unsure{\subset|\in}\widetilde G$ \deleted{trivial, et}\inserted{déjà connu} \deleted{\unsure{l'inverse}}\inserted{et mis} pour mémoire ($\Longrightarrow$ trivial, $\Longleftarrow$ car $\deleted{\widetilde F}F=\mathop{\mathrm{Sup}}\limits_{X\in\widetilde F}X$, $G=\mathop{\mathrm{Sup}}\limits_{Y\in\unsure{\widetilde G|\overline G}}Y$). L'équivalence (i bis) $\Longleftrightarrow$ (ii bis) triviale par l'injectivité de (\deleted{\unsure{*|1.38}}1.39). Or (ii bis) est impliqué par (ii) (tautologique) \unsure{prouvons|poussons} l'inverse, donc que (ii bis) implique que $\mathrm{Mulomb}(F)$ est fermé dans $\mathrm{Mulom}\,G$ pour l'inclusion. Mais \unsure{on aura} $\widetilde F\subset\widetilde G$ \deleted{et il \unsure{est}}\inserted{donc} \deleted{et \unsure{on|en}} $F\leq\unsure{G|C}$, et il est \deleted{\unsure{connu}}\inserted{connu et trivial} que $\widetilde F$ est \inserted{alors} fermé dans $\widetilde G$.

\uline{Corollaire 1} \deleted{Conditions équivalentes\unsure{,} pour}\inserted{Soient} $F,G\in\unsure{M|\mathcal M}$ \inserted{tels que $F\unsure{\between}G$, alors} \deleted{(i) $F\between G$}
\[
 \deleted{\text{(ii)}}\quad \mathrm{Mulomb}(F)\between\mathrm{Mulomb}(G),
 \tag*{(*)}
\]
\deleted{Quand elles sont \unsure{remplies|vérifiées|compatibles},}\inserted{et} on a
\[
 \left\{\begin{aligned}
  \mathrm{Mulomb}(F\vee G)&=\mathrm{Mulom}(F)\cup\mathrm{Mulomb}(G)\\
  \mathrm{Mulomb}(F\wedge G)&=\mathrm{Mulom\unsure{b|s}}(F)\cap\mathrm{Mulomb}(G).
 \end{aligned}\right.
 \tag*{(**)}
\]

\deleted{La condition (ii) signifie que $\widetilde F$ et $\widetilde G$ sont fermés dans $\widetilde F\cup\widetilde G\subset M$, pour la relation $\ll$ (\unsure{induite|induit} par $\subset$ dans $\unsure{\mathcal P|\mathfrak P|\mathscr P|P}(M)$).}
